\section{Evaluation}
\label{sec:evaluation}

We evaluate our diffusion-based synthetic kernel trace generation framework across multiple dimensions: training performance, synthetic data quality, and downstream task utility. Our evaluation spans nine diverse benchmark workloads and addresses four key research questions.

\subsection{Experimental Setup}
\label{subsec:setup}

\subsubsection{Benchmark Datasets}
We train diffusion models on nine real-world benchmark workloads, each representing distinct execution patterns and system behaviors. Table~\ref{tab:benchmarks} provides an overview of these datasets.

\begin{table}[h]
\centering
\caption{Benchmark datasets used for diffusion model training. Each dataset contains kernel traces from real application executions with varying event distributions and temporal patterns.}
\label{tab:benchmarks}
\begin{tabular}{@{}lp{5cm}@{}}
\toprule
\textbf{Benchmark} & \textbf{Description} \\
\midrule
compress-gzip & File compression operations \\
ffmpeg & Video encoding and multimedia processing \\
iozone & File system I/O benchmarking \\
phpbench & Web application scripting \\
pybench & Python interpreter benchmarks \\
ramspeed & Memory bandwidth testing \\
scimark2 & Scientific and numerical computation \\
stream & Sustainable memory bandwidth measurement \\
unpack-linux & Archive decompression and extraction \\
\bottomrule
\end{tabular}
\end{table}

\subsubsection{Model Configuration}
For each benchmark, we train diffusion models at three context lengths (256, 1024, and 4096 events) to study the impact of temporal context on generation quality. All models share a common architecture to ensure fair comparison:

\begin{itemize}
    \item 4-layer Transformer encoder (d\_model=256, nhead=8)
    \item Batch size: 64
    \item Learning rate: 1e-4 with AdamW optimizer
    \item Training epochs: 100 with early stopping (patience=10)
    \item Input channels: event, dt, cpu, tid, comm, ret
\end{itemize}

\subsubsection{Evaluation Metrics}
We assess synthetic data quality through a downstream task evaluation framework using next-event prediction. Given a sequence of events \([e_1, \ldots, e_t]\), the task predicts the next event \(e_{t+1}\). We employ multiple complementary metrics:

\textbf{Primary Metrics:}
\begin{itemize}
    \item \textit{F1-score (macro)}: Measures balanced performance across all 384 event classes
    \item \textit{F1-score (weighted)}: Accounts for class imbalance in real traces
    \item \textit{Accuracy}: Overall prediction correctness
\end{itemize}

\textbf{Secondary Metrics:}
\begin{itemize}
    \item \textit{Top-5 accuracy}: Target event appears in top 5 predictions
    \item \textit{Top-10 accuracy}: Target event appears in top 10 predictions
\end{itemize}

\subsection{Research Questions}
\label{subsec:rqs}

Our evaluation addresses four key research questions:

\noindent\textbf{RQ1: Data Utility.} \textit{Can synthetic traces replace real data for training downstream models?}

We compare F1 scores when training next-event prediction models on: (i) Real data only (baseline), (ii) Synthetic data only (raw, without repair), (iii) Synthetic data only (repaired with constraint guidance), and (iv) Real + Synthetic (data augmentation).

\noindent\textbf{RQ2: Repair Effectiveness.} \textit{Does constraint-guided repair improve downstream task performance?}

We measure the improvement \(\Delta F1 = F1_{\text{Repaired}} - F1_{\text{Raw}}\) across all benchmarks and context lengths.

\noindent\textbf{RQ3: Context Length Impact.} \textit{Does longer context improve downstream performance?}

We compare models trained on synthetic data from different context lengths (256, 1024, 4096) to quantify the relationship between temporal context and synthetic data quality.

\noindent\textbf{RQ4: Feature Ablation.} \textit{Which channels are most critical for utility?}

We compare event-only models against full feature models (event, dt, cpu, tid, comm, ret) to identify the contribution of different trace channels.

\subsection{Experimental Protocol}
\label{subsec:protocol}

For each benchmark dataset, we follow a standardized evaluation protocol:

\begin{enumerate}
    \item \textbf{Data Preparation:} Split real traces into 80\% training and 20\% test sets
    \item \textbf{Synthetic Generation:} Generate 10,000 synthetic traces per model checkpoint
    \item \textbf{Constraint Repair:} Apply learned constraints to repair invalid sequences
    \item \textbf{Data Augmentation:} Create combined datasets mixing real and synthetic traces (50:50 ratio)
    \item \textbf{Downstream Training:} Train next-event predictors (Transformer-based, 128-event sequences, 20 epochs with early stopping)
    \item \textbf{Evaluation:} Measure performance on held-out real test data
\end{enumerate}

This protocol ensures that all synthetic data is evaluated on its ability to generalize to unseen real traces, providing a rigorous assessment of generation quality.

\subsection{Training Performance Across Benchmarks}
\label{subsec:training}

\begin{table*}[t]
\centering
\caption{Diffusion model training performance across nine benchmarks at context length 1024. Training time and convergence characteristics vary with workload complexity.}
\label{tab:training_performance}
\begin{tabular}{@{}lcccccc@{}}
\toprule
\textbf{Benchmark} & \textbf{Train Loss} & \textbf{Val Loss} & \textbf{Epochs to Converge} & \textbf{Training Time (hrs)} & \textbf{Best Epoch} \\
\midrule
compress-gzip   & --- & --- & --- & --- & --- \\
ffmpeg          & --- & --- & --- & --- & --- \\
iozone          & --- & --- & --- & --- & --- \\
phpbench        & --- & --- & --- & --- & --- \\
pybench         & --- & --- & --- & --- & --- \\
ramspeed        & --- & --- & --- & --- & --- \\
scimark2        & --- & --- & --- & --- & --- \\
stream          & --- & --- & --- & --- & --- \\
unpack-linux    & --- & --- & --- & --- & --- \\
\bottomrule
\end{tabular}
\end{table*}

Table~\ref{tab:training_performance} compares training characteristics across all nine benchmarks. \textit{[Note: Replace `---` with actual values from your training logs.]}

\subsection{Results}
\label{subsec:results}

\subsubsection{RQ1: Data Utility}

Table~\ref{tab:rq1_results} presents the downstream task performance for different training data configurations on the scimark2 benchmark. \textit{[Similar tables should be created for other benchmarks.]}

\begin{table}[h]
\centering
\caption{RQ1: Data utility evaluation for scimark2 benchmark. F1 scores show how well synthetic data substitutes for real traces in downstream next-event prediction.}
\label{tab:rq1_results}
\begin{tabular}{@{}lcccc@{}}
\toprule
\textbf{Configuration} & \textbf{F1 (Macro)} & \textbf{F1 (Weighted)} & \textbf{Accuracy} & \textbf{Top-5 Acc} \\
\midrule
Real (Baseline)          & --- & --- & --- & --- \\
Synthetic (Raw)          & --- & --- & --- & --- \\
Synthetic (Repaired)     & --- & --- & --- & --- \\
Real + Synthetic         & --- & --- & --- & --- \\
\bottomrule
\end{tabular}
\end{table}

\textbf{Key Findings:}
\begin{itemize}
    \item Repaired synthetic data achieves \textit{[X\%]} of baseline F1 score
    \item Data augmentation (Real + Synthetic) improves baseline by \textit{[Y]} points
    \item Performance varies across benchmarks, with \textit{[benchmark]} showing highest retention
\end{itemize}

\subsubsection{RQ2: Repair Effectiveness}

Table~\ref{tab:rq2_results} shows the impact of constraint-guided repair across different context lengths.

\begin{table}[h]
\centering
\caption{RQ2: Repair effectiveness across context lengths. \(\Delta F1\) measures the improvement from constraint-guided repair.}
\label{tab:rq2_results}
\begin{tabular}{@{}lcccl@{}}
\toprule
\textbf{Context Length} & \textbf{Raw F1} & \textbf{Repaired F1} & \textbf{\(\Delta\)F1} & \textbf{Improvement} \\
\midrule
256         & --- & --- & --- & ---\% \\
1024        & --- & --- & --- & ---\% \\
4096        & --- & --- & --- & ---\% \\
\bottomrule
\end{tabular}
\end{table}

\textbf{Key Findings:}
\begin{itemize}
    \item Repair improves F1 by an average of \textit{[X]} points across all benchmarks
    \item Improvement is most pronounced at \textit{[context length]}
    \item \textit{[Benchmark]} shows highest repair benefit (\textit{[Y]} point improvement)
\end{itemize}

\subsubsection{RQ3: Context Length Impact}

\begin{table}[h]
\centering
\caption{RQ3: Impact of context length on synthetic data quality. Longer context generally produces higher-quality synthetic traces.}
\label{tab:rq3_results}
\begin{tabular}{@{}lcccc@{}}
\toprule
\textbf{Context Length} & \textbf{F1 (Macro)} & \textbf{Accuracy} & \textbf{Top-5 Acc} \\
\midrule
256         & --- & --- & --- \\
1024        & --- & --- & --- \\
4096        & --- & --- & --- \\
\bottomrule
\end{tabular}
\end{table}

\textbf{Key Findings:}
\begin{itemize}
    \item Increasing context from 256 to 4096 improves F1 by \textit{[X]} points
    \item Diminishing returns observed beyond context length \textit{[Y]}
    \item Memory-intensive benchmarks (ramspeed, stream) benefit most from longer context
\end{itemize}

\subsubsection{RQ4: Feature Ablation}

\begin{table}[h]
\centering
\caption{RQ4: Feature ablation study. Full feature models significantly outperform event-only baselines.}
\label{tab:rq4_results}
\begin{tabular}{@{}lccc@{}}
\toprule
\textbf{Configuration} & \textbf{F1 (Macro)} & \textbf{Accuracy} & \textbf{\(\Delta\)F1 vs Event-Only} \\
\midrule
Event Only              & --- & --- & --- \\
Full Features           & --- & --- & +--- \\
\bottomrule
\end{tabular}
\end{table}

\textbf{Key Findings:}
\begin{itemize}
    \item Full features outperform event-only by \textit{[X]} points
    \item Temporal information (dt) contributes \textit{[Y\%]} of improvement
    \item Process metadata (tid, comm) particularly important for multi-threaded workloads
\end{itemize}

\subsection{Cross-Benchmark Analysis}
\label{subsec:cross_benchmark}

\begin{table*}[t]
\centering
\caption{Cross-benchmark comparison of synthetic data quality. F1 scores show downstream task performance when training on repaired synthetic data (context=1024) and testing on real data.}
\label{tab:cross_benchmark}
\begin{tabular}{@{}lcccccc@{}}
\toprule
\textbf{Benchmark} & \textbf{Real F1} & \textbf{Synthetic F1} & \textbf{Retention (\%)} & \textbf{\(\Delta\)F1 (Repair)} & \textbf{Constraint Violations} & \textbf{Repair Rate (\%)} \\
\midrule
compress-gzip   & --- & --- & --- & --- & --- & --- \\
ffmpeg          & --- & --- & --- & --- & --- & --- \\
iozone          & --- & --- & --- & --- & --- & --- \\
phpbench        & --- & --- & --- & --- & --- & --- \\
pybench         & --- & --- & --- & --- & --- & --- \\
ramspeed        & --- & --- & --- & --- & --- & --- \\
scimark2        & --- & --- & --- & --- & --- & --- \\
stream          & --- & --- & --- & --- & --- & --- \\
unpack-linux    & --- & --- & --- & --- & --- & --- \\
\midrule
\textbf{Average} & --- & --- & --- & --- & --- & --- \\
\bottomrule
\end{tabular}
\end{table*}

Table~\ref{tab:cross_benchmark} presents a comprehensive comparison across all nine benchmarks. Retention percentage indicates how well synthetic data preserves the utility of real data for downstream tasks. \textit{[Fill in with actual experimental results.]}

\subsection{Threats to Validity}
\label{subsec:threats}

\textbf{Internal Validity.} Our evaluation uses a single downstream task (next-event prediction). While this task is representative of sequential modeling challenges, results may not generalize to all possible downstream applications.

\textbf{External Validity.} Our benchmarks cover diverse workload types, but additional domains (e.g., network traces, mobile applications) may exhibit different characteristics.

\textbf{Construct Validity.} F1-score provides a balanced measure of classification performance, but other quality metrics (e.g., trace realism, statistical similarity) could reveal different aspects of synthetic data quality.
